\chapter{Diseases of the Lymphatic and Immune System}
\label{ch:lymph}
The lymphatic system includes lymph nodes, lymphatic vessels, tonsils, spleen, thymus, and bone marrow. It functions in fluid homeostasis, immune surveillance, and fat absorption. The immune system encompasses innate and adaptive immunity, protecting against infections and malignancy. Disorders include autoimmune diseases, immunodeficiencies, and hypersensitivity reactions.

%-------------------------------------------------------------
\section{Autoimmune Diseases}
%-------------------------------------------------------------

%-------------------------------------------------------------

%-------------------------------------------------------------

%-------------------------------------------------------------

%-------------------------------------------------------------

%-------------------------------------------------------------

\label{sec:Autoimmune Diseases}
%-------------------------------------------------------------%-------------------------------------------------------------

\subsection{Antiphospholipid Syndrome}
\label{sec:Antiphospholipid Syndrome}
Antiphospholipid syndrome is an autoimmune thrombophilic disorder characterized by arterial or venous thrombosis and/or pregnancy morbidity in the presence of persistent antiphospholipid antibodies: anticardiolipin antibodies, anti-$\beta$2-glycoprotein I antibodies, and/or lupus anticoagulant, and may be primary (no associated disease) or secondary (associated with SLE). Clinical features include deep vein thrombosis, pulmonary embolism, stroke, myocardial infarction, recurrent pregnancy loss, livedo reticularis, thrombocytopenia, and cardiac valve disease, with catastrophic APS (CAPS) being a rare, life-threatening variant with widespread thrombosis affecting multiple organs simultaneously. Diagnosis requires at least one clinical criterion (thrombosis or pregnancy morbidity) and one laboratory criterion (persistent antiphospholipid antibodies on two occasions $\geq$12 weeks apart). Management is long-term anticoagulation with warfarin (INR target 2--3 for venous thrombosis, 2--3 or 3--4 for arterial thrombosis), with DOACs not recommended for APS. Prognosis depends on the site and severity of thrombosis.

\subsection{Behçet's Disease}
\label{sec:Behçet's Disease}
Behçet's disease is a systemic vasculitis of unknown etiology, characterized by recurrent oral aphthous ulcers (the hallmark feature, present in $>$98\% of patients), genital ulcers, uveitis (which can be sight-threatening), and skin lesions (erythema nodosum, pseudofolliculitis, pathergy), with a predilection for the Silk Road populations (Turkey, Iran, Japan, China), with a prevalence of up to 1 in 250 in Turkey. The condition results from systemic vasculitis. Clinical features also include vascular involvement (arterial aneurysms, venous thrombosis), neurological involvement (neuro-Behçet's with brainstem/diencephalic lesions), gastrointestinal involvement (ileocecal ulcers), and arthritic involvement. Diagnosis is clinical, based on the International Criteria for Behçet's Disease (ICBD), with no pathognomonic laboratory test. Management is tailored to organ involvement: topical agents for oral/genital ulcers, colchicine for mucocutaneous disease, azathioprine or TNF inhibitors for ocular or systemic involvement, and cyclophosphamide for severe vascular or neurological disease. Prognosis is variable depending on organ involvement.

\subsection{Rheumatoid Arthritis}
\label{sec:Rheumatoid Arthritis (Immune)}
Rheumatoid arthritis is discussed in detail in Chapter~\ref{ch:bones}. It is a systemic autoimmune disease primarily affecting joints but with significant extrarticular manifestations involving the immune system. Autoantibodies (rheumatoid factor, anti-CCP) are central to diagnosis and prognosis.

\subsection{Sarcoidosis}
\label{sec:lymph:Sarcoidosis}
Sarcoidosis is discussed in detail in Chapter~\ref{ch:lungs}. It is a multisystem granulomatous disease of unknown etiology affecting the lungs, lymph nodes, skin, eyes, and other organs. Non-caseating granulomas on biopsy are the histological hallmark.

\subsection{Sjögren's Syndrome}
\label{sec:Sjögren's Syndrome}
Sjögren's syndrome is a chronic autoimmune disease characterized by lymphocytic infiltration of exocrine glands (salivary and lacrimal), resulting in dry eyes (keratoconjunctivitis sicca) and dry mouth (xerostomia), predominantly affecting women (9:1), and may be primary (isolated) or secondary (associated with another autoimmune disease, most commonly RA). Clinical features include Sicca symptoms (dry eyes, dry mouth), parotid gland enlargement, dental caries, dysphagia, and extraglandular manifestations (arthritis, interstitial lung disease, renal tubular acidosis, peripheral neuropathy, lymphoma). Diagnosis requires demonstration of glandular inflammation: salivary gland biopsy showing focal lymphocytic sialadenitis (focus score $\geq$1), positive anti-SSA/Ro and anti-SSB/La antibodies, abnormal ocular staining (Schirmer's test, Rose Bengal), and unstimulated salivary flow. Management is symptomatic: artificial tears, pilocarpine or cevimeline (muscarinic agonists to stimulate salivary and lacrimal secretion), and immunosuppression for severe extraglandular manifestations. Prognosis is generally favorable, though patients have a 5--10\% lifetime risk of developing B-cell lymphoma (MALT lymphoma of the parotid gland most common).

\subsection{Systemic Lupus Erythematosus}
\label{sec:Systemic Lupus Erythematosus}
Systemic lupus erythematosus is a chronic, multisystem autoimmune disease characterized by loss of tolerance to self-antigens and production of a wide array of autoantibodies, leading to immune complex deposition and inflammation in multiple organs, predominantly affecting women of childbearing age (9:1 female predominance) and being more common in African Americans, Hispanics, and Asians. The condition results from autoantibody production and immune complex deposition. Clinical features are protean and include malar rash, photosensitivity, oral ulcers, arthritis, serositis (pleuritis, pericarditis), nephritis (lupus nephritis---the most serious manifestation), cytopenias, seizures, and psychosis. The 2019 EULAR/ACR classification criteria require an ANA titer $\geq$1:80 plus weighted criteria across clinical and immunological domains. Diagnosis involves antinuclear antibody (ANA, sensitive but not specific), anti-dsDNA (specific, correlates with nephritis), anti-Smith (highly specific), complement levels (C3, C4 decreased during flares), and renal biopsy for lupus nephritis classification (ISN/RPS). Management includes hydroxychloroquine (cornerstone for all patients, reduces flares and mortality), corticosteroids, and immunosuppressive agents (mycophenolate, azathioprine, cyclophosphamide for severe organ involvement), with newer targeted therapies including belimumab (anti-BLys/BAFF monoclonal antibody), voclosporin, and anifrolumab (anti-type I interferon receptor). Prognosis is variable, with improved outcomes due to early diagnosis and better management.

%-------------------------------------------------------------
\section{Chronic Fatigue Syndrome / Myalgic Encephalomyelitis}
%-------------------------------------------------------------

%-------------------------------------------------------------

%-------------------------------------------------------------

%-------------------------------------------------------------

%-------------------------------------------------------------

%-------------------------------------------------------------

\label{sec:Chronic Fatigue Syndrome / Myalgic Encephalomyelitis}
%-------------------------------------------------------------%-------------------------------------------------------------

\subsection{Chronic Fatigue Syndrome}
\label{sec:Chronic Fatigue Syndrome}
Chronic fatigue syndrome/myalgic encephalomyelitis (CFS/ME) is a complex, debilitating multisystem illness affecting 0.2--0.4\% of the population, with a female predominance. The etiology is unknown but is likely multifactorial, involving post-infectious triggers (EBV, enterovirus, Q fever, Ross River virus), immune dysregulation, neuroendocrine dysfunction, and genetic predisposition, with proposed mechanisms including chronic immune activation, hypothalamic-pituitary-adrenal axis hypofunction, mitochondrial dysfunction, and central sensitization. Clinical features (CDC/Institute of Medicine 2015 diagnostic criteria) include substantial reduction in previous activity level persisting $>$6 months, post-exertional malaise (PEM), unrefreshing sleep, cognitive impairment (brain fog), and orthostatic intolerance, along with additional symptoms such as chronic pain, sore throat, tender lymphadenopathy, and flu-like malaise. Diagnosis is clinical (no biomarkers), requiring ruling out other causes. Management includes symptom management with pacing/activity management (avoiding PEM), and symptom-specific treatments such as sleep hygiene, low-dose naltrexone, and SSRIs. Prognosis is variable, with approximately 5\% recovering fully and most patients experiencing a fluctuating course.

%-------------------------------------------------------------
\section{Hypersensitivity Disorders}
%-------------------------------------------------------------

%-------------------------------------------------------------

%-------------------------------------------------------------

%-------------------------------------------------------------

%-------------------------------------------------------------

%-------------------------------------------------------------

\label{sec:Hypersensitivity Disorders}
%-------------------------------------------------------------%-------------------------------------------------------------

\subsection{Anaphylaxis}
\label{sec:Anaphylaxis}
Anaphylaxis is a severe, life-threatening systemic hypersensitivity reaction (usually IgE-mediated) that occurs rapidly and can involve multiple organ systems. The condition results from exposure to triggers including foods (nuts, shellfish, milk, eggs), insect stings (Hymenoptera), medications (penicillin, NSAIDs), latex, and contrast media. Clinical features include skin involvement (urticaria, angioedema, flushing in 80--90\%), respiratory compromise (laryngeal edema, bronchospasm), cardiovascular collapse (hypotension, distributive shock), gastrointestinal symptoms (nausea, vomiting, diarrhea), and neurological symptoms (anxiety, altered consciousness). Diagnosis is clinical: acute onset (minutes to hours) of skin/mucosal involvement plus respiratory compromise or hypotension after exposure to a likely allergen, with serum tryptase possibly elevated but not required for diagnosis. Management is with intramuscular epinephrine (0.3--0.5 mg of 1:1000 in the mid-ant thigh, repeat every 5--15 minutes as needed), IV fluid resuscitation, H1 and H2 antihistamines, corticosteroids, and inhaled bronchodilators for bronchospasm. Prognosis is excellent with prompt epinephrine administration.

\subsection{Drug Hypersensitivity}
\label{sec:Drug Hypersensitivity}
Drug hypersensitivity reactions are adverse drug reactions with an immunologic basis, representing approximately 5--10\% of all adverse drug reactions. The condition may be immediate (IgE-mediated, occurring within 1--6 hours: urticaria, angioedema, anaphylaxis) or delayed (T-cell mediated, occurring hours to days after exposure: maculopapular exanthema, DRESS syndrome, Stevens-Johnson syndrome, toxic epidermal necrolysis, acute generalized exanthematous pustulosis). Clinical features depend on the type and severity of reaction. Common culprits include beta-lactam antibiotics, sulfonamides, anticonvulsants, and NSAIDs. Diagnosis involves detailed clinical history, skin testing (for IgE-mediated reactions), drug provocation tests, and lymphocyte transformation tests (for delayed reactions). Management includes drug avoidance, acute treatment of reactions, and in some cases, desensitization protocols for essential medications. Prognosis is generally favorable with avoidance of the offending drug.

\subsection{Stevens-Johnson Syndrome and Toxic Epidermal Necrolysis}
\label{sec:Stevens-Johnson Syndrome and Toxic Epidermal Necrolysis}
Stevens-Johnson syndrome and toxic epidermal necrolysis are severe, life-threatening mucocutaneous reactions characterized by extensive keratinocyte apoptosis and epidermal detachment, with SJS affecting $<$10\% of body surface area, SJS/TEN overlap affecting 10--30\%, and TEN affecting $>$30\% BSA. The most common cause is drug hypersensitivity (anticonvulsants [carbamazepine, lamotrigine, phenytoin], allopurinol, sulfonamides, NSAIDs, antibiotics). Clinical features include prodromal fever and malaise, followed by painful, dusky red to purpuric macules that coalesce and progress to epidermal detachment with flaccid blisters and erosion, with mucosal involvement (oral, ocular, genital) present in over 90\%. Diagnosis is clinical. Management requires burn unit or ICU care: fluid resuscitation, wound care, nutritional support, ophthalmological consultation, pain management, and withdrawal of the offending drug, with cyclosporine and IVIG having shown some benefit. Prognosis is guarded, with mortality of 10--30\% for SJS and up to 50\% for TEN.

%-------------------------------------------------------------
\section{Immunodeficiency Disorders}
%-------------------------------------------------------------

%-------------------------------------------------------------

%-------------------------------------------------------------

%-------------------------------------------------------------

%-------------------------------------------------------------

%-------------------------------------------------------------

\label{sec:Immunodeficiency Disorders}
%-------------------------------------------------------------%-------------------------------------------------------------

\subsection{Common Variable Immunodeficiency}
\label{sec:Common Variable Immunodeficiency}
Common variable immunodeficiency is the most common symptomatic primary antibody deficiency, affecting approximately 1 in 25,000--50,000. The condition is characterized by hypogammaglobulinemia (decreased IgG and IgA, with or without decreased IgM), impaired vaccine antibody responses, and increased susceptibility to bacterial infections, particularly sinopulmonary infections (pneumonia, sinusitis, bronchitis). Diagnosis requires exclusion of other causes of hypogammaglobulinemia and is typically made in the second to fourth decade of life. Clinical features include recurrent infections. Complications include bronchiectasis (from recurrent infections), autoimmune cytopenias, granulomatous disease, lymphoproliferation, and increased risk of gastrointestinal and lymphoid malignancies. Management is with immunoglobulin replacement therapy (intravenous or subcutaneous), prophylactic antibiotics, and management of complications. Prognosis is variable but improved with immunoglobulin replacement.

\subsection{DiGeorge Syndrome}
\label{sec:DiGeorge Syndrome}
DiGeorge syndrome (22q11.2 deletion syndrome, velocardiofacial syndrome) is a developmental disorder caused by a microdeletion on chromosome 22q11.2, resulting in developmental abnormalities derived from the third and fourth pharyngeal arches. Clinical features include thymic hypoplasia or aplasia (T-cell immunodeficiency of varying severity), hypoparathyroidism (hypocalcemia), congenital heart defects (tetralogy of Fallot, interrupted aortic arch, truncus arteriosus, ventricular septal defect), characteristic facies (broad nasal bridge, small ears, micrognathia), palatal abnormalities (velopharyngeal insufficiency, cleft palate), and learning difficulties, with the severity of immunodeficiency varying widely. Diagnosis is confirmed by FISH or chromosomal microarray showing 22q11.2 deletion. Management addresses individual manifestations: cardiac surgery, calcium/vitamin D supplementation for hypocalcemia, thymic transplantation for complete aplasia, and developmental/educational support. Prognosis depends on the severity of associated anomalies.

\subsection{Hereditary Angioedema}
\label{sec:Hereditary Angioedema}
Hereditary angioedema (HAE) is an autosomal dominant disorder caused by deficiency (type I, 85\%) or dysfunction (type II, 15\%) of C1 esterase inhibitor (C1-INH), with a prevalence of 1 in 50,000. The condition results from uncontrolled activation of the contact system, leading to excessive bradykinin production, increased vascular permeability, and recurrent episodes of subcutaneous and submucosal edema. Clinical features include recurrent episodes of non-pitting, non-pruritic edema affecting the extremities, face, genitals, trunk, and potentially life-threatening laryngeal edema leading to airway obstruction. Gastrointestinal involvement causes severe abdominal pain, nausea, vomiting, and diarrhea, often misdiagnosed as acute abdomen. Episodes last 2--5 days without treatment. Unlike allergic angioedema, there is no urticaria, pruritus, or response to antihistamines or epinephrine. Diagnosis is based on low C4 levels during attacks (screening test), low C1-INH antigenic levels (type I), low functional C1-INH activity (type II), and C1-INH genetic testing. Management includes on-demand treatment for acute attacks (plasma-derived C1-INH, recombinant C1-INH, icatibant [bradykinin B2 receptor antagonist], ecallantide [kallikrein inhibitor]), short-term prophylaxis prior to procedures (plasma-derived C1-INH or attenuated androgens), and long-term prophylaxis (plasma-derived C1-INH, lanadelumab [monoclonal antibody against plasma kallikrein], berotralstat [oral plasma kallikrein inhibitor], or attenuated androgens: danazol, stanozolol). Prognosis is excellent with appropriate management; before modern therapy, mortality from laryngeal edema was 25--30\%.

\subsection{HIV/AIDS}
\label{sec:HIV/AIDS}
Human immunodeficiency virus (HIV) infection is a progressive immune deficiency caused by HIV-1 (global pandemic) or HIV-2 (West Africa), ultimately leading to acquired immunodeficiency syndrome (AIDS) if untreated. The condition results from targeting of CD4+ T lymphocytes, macrophages, and dendritic cells, with transmission occurring through sexual contact, blood-borne exposure, and vertical transmission (mother to child). Clinical features follow a natural history including acute retroviral syndrome (fever, rash, lymphadenopathy, myalgia 2--4 weeks after infection), clinical latency (chronic infection with asymptomatic period lasting years), and AIDS (CD4 count $<$200 cells/$\mu$L or presence of AIDS-defining opportunistic infections and malignancies). Diagnosis is by fourth-generation HIV antigen/antibody combination immunoassay, with confirmatory testing using HIV-1/HIV-2 differentiation immunoassay and HIV RNA viral load. Management is with antiretroviral therapy (ART), typically a combination of three or more drugs from different classes (NRTIs, NNRTIs, integrase inhibitors, protease inhibitors, entry inhibitors), with goals of viral suppression (undetectable viral load), immune reconstitution (CD4 recovery), and prevention of transmission (U=U: Undetectable = Untransmittable). Prognosis has been transformed by ART, with near-normal life expectancy for adherent patients.

\subsection{Severe Combined Immunodeficiency}
\label{sec:Severe Combined Immunodeficiency}
Severe combined immunodeficiency is a group of rare, life-threatening genetic disorders characterized by severe impairment of both cellular (T-cell) and humoral (B-cell) immunity, with or without natural killer (NK) cell deficiency. The condition results from over 20 genetic causes, with X-linked SCID (IL2RG mutations, most common) and adenosine deaminase (ADA) deficiency being the best known; without treatment, affected infants typically die within the first year of life from severe, recurrent, opportunistic infections. Clinical features include recurrent severe infections, failure to thrive, and persistent diarrhea. Newborn screening detects SCID by measuring T-cell receptor excision circles (TRECs) on dried blood spots, with diagnosis confirmed by lymphocyte subset analysis (flow cytometry), immunoglobulin levels, and genetic testing. Management includes hematopoietic stem cell transplantation (curative, ideally before 3.5 months of age), enzyme replacement therapy (for ADA deficiency), and gene therapy (for X-linked SCID and ADA-SCID). Prognosis is excellent with early transplantation.

%-------------------------------------------------------------
\section{Infectious Mononucleosis}
%-------------------------------------------------------------

%-------------------------------------------------------------

%-------------------------------------------------------------

%-------------------------------------------------------------

%-------------------------------------------------------------

%-------------------------------------------------------------

\label{sec:section:Infectious Mononucleosis}
%-------------------------------------------------------------%-------------------------------------------------------------

\subsection{Infectious Mononucleosis}
\label{sec:Infectious Mononucleosis}
Infectious mononucleosis is an acute lymphoproliferative disorder caused by Epstein-Barr virus (EBV, human herpesvirus 4), primarily affecting adolescents and young adults. The condition results from EBV infecting B lymphocytes, triggering a robust CD8+ T-cell response that produces the characteristic atypical lymphocytosis and many of the clinical manifestations, with transmission through saliva (``kissing disease'') and an incubation period of 4--6 weeks. Clinical features include the classic triad of fever, pharyngitis (with severe exudative tonsillitis), and lymphadenopathy (typically posterior cervical), along with fatigue (often profound and prolonged), hepatosplenomegaly, maculopapular rash (10--15\%, but up to 90\% if amoxicillin or ampicillin is given), periorbital edema, and palatal petechiae. Diagnosis is suggested by lymphocytosis ($>$50\% lymphocytes, with $>$10\% atypical lymphocytes), positive heterophile antibody test (Monospot), and EBV-specific serology (anti-VCA IgM for acute infection, anti-VCA IgG for past infection, anti-EBNA for past infection). Management is symptomatic: rest, hydration, analgesics, and antipyretics, with corticosteroids considered for impending airway obstruction, severe thrombocytopenia, or hemolytic anemia. Prognosis is generally excellent, though splenic rupture is a rare but life-threatening complication.

